% =============================================================================
% 00_ABSTRACT — Literature Review Abstract
% =============================================================================
% Word count target: ~300 words
% Structure: Context → Objective → Method → Findings → Implications
% =============================================================================

\begin{abstract}
% -----------------------------------------------------------------------------
% CONTEXT — Why does this matter? Set the societal/stakes context.
% -----------------------------------------------------------------------------
% [TODO: 2–3 sentences]
% • AI is rapidly deployed in clinical settings worldwide.
% • Unlike drugs, AI systems evolve post-deployment → continuous performance
%   governance is structurally absent from traditional regulatory pathways.
% • Gaps between pilot studies and routine care have been widely documented
%   but insufficiently synthesized from a governance-and-evaluation angle.
%
The integration of artificial intelligence (AI) into healthcare systems has
accelerated dramatically over the past five years, with hundreds of AI-enabled
medical devices authorized by regulators such as the FDA and dozens of large
language models entering clinical pilots. Unlike pharmaceutical products, AI
systems are not static: they evolve through continuous learning, drift over
time, and may perform differently under real-world conditions than in controlled
validation studies. This poses a fundamental challenge to traditional
regulatory frameworks designed for static medical devices.

% -----------------------------------------------------------------------------
% OBJECTIVE — What is the aim of this review? State the research question.
% -----------------------------------------------------------------------------
% [TODO: 1–2 sentences]
% • State the specific question this review addresses.
% • Mention the thematic lens (evaluation methods + governance structures).
% • Mention scope (real-world deployment, health AI broadly).
%
This literature review systematically examines the question of \textbf{how
health AI systems are evaluated and governed in real-world deployment contexts}.
Specifically, it asks: (\textit{i}) what evaluation frameworks and methods are
used to assess AI performance, reliability, and safety post-deployment;
(\textit{ii}) what governance structures shape the sustainable implementation
of clinical AI; and (\textit{iii}) what gaps remain at the intersection of
technical evaluation and societal impact.

% -----------------------------------------------------------------------------
% METHOD — How was the review conducted?
% -----------------------------------------------------------------------------
% [TODO: 2–3 sentences]
% • Search strategy: which databases, time window, keywords.
% • Selection criteria: inclusion/exclusion (peer-reviewed, English, health AI).
% • Thematic synthesis approach.
%
A systematic search was conducted across PubMed, Semantic Scholar, arXiv, and
DBLP for peer-reviewed publications and preprints published between 2020 and 2026,
using a curated keyword set spanning five thematic axes: AI Evaluation and
Validation, Participatory Governance, Adaptive Regulation, Evidence and
Implementation, and Clinical Health AI. A total of approximately 2,700 papers
were identified across sources, reduced to a final corpus of 311 papers after
deduplication and quality screening. Of these, 78 (25.1\%) were designated
Priority 3 (composite score ≥ 7.0) and form the primary evidence base; the
remainder were retained as supporting literature. A thematic synthesis was
performed, organizing findings across three analytical dimensions: (1) evaluation
methods and metrics, (2) governance and regulatory frameworks, and (3)
implementation outcomes and barriers.

% -----------------------------------------------------------------------------
% FINDINGS — What did we learn? Key themes (3–4 sentences).
% -----------------------------------------------------------------------------
% [TODO: 3–4 sentences covering main findings]
% • Theme 1: Evaluation gap between lab/real-world.
% • Theme 2: Governance frameworks emerging (FDA, EU, multi-stakeholder).
% • Theme 3: Pilot trap phenomenon.
% • Theme 4: Ethical concerns and participatory governance.
%
The review identifies three major structural gaps in the health AI literature.
First, a pervasive \textbf{ecological validity gap}: the majority of
validation studies are conducted under controlled conditions and fail to capture
the temporal dynamics, contextual reasoning, and multi-professional workflows
of real clinical environments \citep{turchi2025ecological}. Second, a
\textbf{governance capacity gap}: existing regulatory frameworks (FDA Software
Pre-Cert, EU AI Act) specify \textit{what} governance should achieve but
assume pre-existing organizational capacity that health systems have not yet
developed \citep{tian2026pilottrap}. Third, an \textbf{accountability gap}:
AI systems operating under continuous learning protocols may drift beyond their
original validation parameters without mandatory post-market performance
monitoring, creating untracked risk for patient safety \citep{abulibdeh2025illusion}.

% -----------------------------------------------------------------------------
% IMPLICATIONS — So what? Why does this matter for the field?
% -----------------------------------------------------------------------------
% [TODO: 2–3 sentences]
% • Implications for regulators, implementers, clinicians.
% • What a governance-and-evaluation integrated framework would require.
% • Open research questions.
%
These findings imply that achieving safe and equitable health AI deployment
requires an \textbf{integrated governance-and-evaluation approach} that moves
beyond pre-market authorization alone. Regulators must mandate continuous
post-market performance monitoring with standardized metrics; health systems
must invest in institutional capacity-building as a precondition for sustainable
AI integration, not an afterthought; and participatory mechanisms must be
institutionalized to ensure that AI systems reflect the needs and values of the
communities they serve. The review concludes with a research agenda outlining
the five most pressing unanswered questions at the intersection of health AI
evaluation and governance.

% =============================================================================
% KEYWORDS (add after abstract body)
% =============================================================================
\medskip

\noindent\textbf{Keywords:} health AI, AI evaluation, algorithmic governance,
clinical AI deployment, adaptive regulation, participatory governance,
implementation science, post-market surveillance, AI validation, digital health

\end{abstract}